% \iffalse
%<*driver>
\ProvidesFile{seqenseigne-core-exos.dtx}%
[2026/05/04 v1.0-dev seqenseigne – Module exercices et remédiation]
%</driver>
% \fi
%
% =============================================================================
% seqenseigne-core-exos.dtx
% Module : CORE-EXOS
% Rôle   :
%   - Regrouper toutes les macros et environnements liés aux EXERCICES
%     d'un livret de séquence (extrait du core en v0.11.5).
%   - Fournir le mécanisme de REMÉDIATION : pour chaque exercice de
%     série fondamentale (typiquement F), permettre d'attacher un énoncé
%     et un corrigé de remédiation, qui sortiront dans deux sections
%     distinctes en fin de livret.
%   - Fournir un CADRE DE RÉPONSE prêt à l'emploi pour qu'un élève
%     puisse répondre directement dans le livret imprimé.
%
% Dépendances :
%   - seqenseigne-core.dtx  (compteurs ExoNum/SerieExosNum, environnements
%                             boiteTitreExo, boitePaleBreakNoTitle,
%                             macros \seqTitreSection, \seqTitreSectionCorrige)
%   - seqenseigne-theme.dtx  (boiteVariableTitre pour le cadre de réponse)
%   - paquet 'multicol', 'tcolorbox', 'etoolbox', 'ifthen' (chargés par core)
%
% Macros publiques :
%   - \begin{seqSerieExos}{N} ... \end{seqSerieExos}%
%   - \begin{seqExercice}[options] ... \end{seqExercice}%
%   - \seqCorrige{texte}%
%   - \seqRemediation{énoncé}{corrigé}     (v0.11.5)
%   - \seqCadreReponse{N}                   (v0.11.5)
%   - \seqInitCorriges
%   - \seqInitRemediation                   (v0.11.5)
%   - \seqAfficheCorriges
%   - \seqAfficheRemediations               (v0.11.5)
%   - \seqAfficheCorrigesRemediation        (v0.11.5)
%   - \begin{seqQcm}[options] ... \end{seqQcm} (v0.13.5.2 étape 2,
%       refactorisé en environnement enfant utilisable dans
%       seqEvalExercice ou seqExercice)
%   - \seqQcmQuestion{texte question} (v0.13.5.2 étape 2, produit le
%       numéro de la question (calculé via \therownum) et son texte ;
%       les & et \\ doivent être écrits en clair par l'utilisateur)
%
% v0.13.5.2 étape 1 — Déplacement depuis seqenseigne-core.dtx :
%   - environnement qcm (un QCM est un format particulier d'exercice,
%     utilisable en évaluation comme en livret de séquence ; sa place
%     est ici plutôt qu'avec les macros d'évaluation)
%   - helpers QCM partagés
%
% v0.13.5.2 étape 2 — Renommages et refactorisation :
%   - environnement qcm                  → seqQcm (environnement enfant)
%   - macro \seqEvalQCMItem              → \seqQcmQuestion (cf. ci-dessous)
%   - compteur i                         → seq@qcm@iter (préfixé interne)
%   - compteur AlphCode                  → seq@qcm@alphcode (préfixé interne)
%   - commande \repCodes                 → \seq@qcm@repCodes (préfixée interne)
%   - compteur QcmItemNumCnt             → SUPPRIMÉ (cf. ci-dessous)
%
% L'environnement seqQcm ne gère plus le cartouche d'exercice, le numéro
% d'exercice, ni l'écriture du corrigé différé. Ces responsabilités sont
% reportées sur l'environnement parent (seqEvalExercice en éval, ou
% seqExercice en livret de séquence).
%
% Le scoring (paramètres ptOK/ptPartiel/ptKO) est géré en tout-ou-rien :
% les règles de scoring ne sont affichées que si les TROIS paramètres
% sont renseignés. Cas d'usage : un QCM AE en révision dans un livret
% de séquence n'a aucun des trois renseigné, pas de scoring affiché.
%
% Numérotation des questions du QCM :
%   - À l'origine, un compteur \seqQcmItemNum était utilisé via une
%     macro \seqQcmItem qui faisait \stepcounter à chaque appel.
%   - Mais tabularray fait plusieurs passes pour stabiliser sa mise en
%     page, ce qui faisait sauter la numérotation (1, 2 puis 4, 5...).
%   - Solution adoptée : utiliser \therownum (compteur de ligne fourni
%     par tabularray, robuste face aux passes multiples), avec un
%     décalage de -2 pour les 2 lignes d'en-tête.
%   - La macro \seqQcmQuestion fournit ce numéro et le texte de la
%     question. Les & séparateurs et le \\ final doivent être écrits
%     en clair par l'utilisateur (tabularray refuse les macros qui
%     contiennent ces éléments structurels).
%
% Limites tabularray sur seqQcm :
%   - colspec figé en \"X[5,c,m]*{5}{X[2,c,m]}\" (pas de macro à
%     l'intérieur de la spécification de colonnes). D'où la limite à
%     5 colonnes de réponses maximum (mais 2 à 4 marchent visuellement
%     même si on a réservé 5 colonnes).
%   - \cmidrule{2-5} en dur dans l'en-tête (même raison). Le trait
%     s'adapte visuellement même avec moins de réponses effectives.
%
% Mécanisme de remédiation :
%   La remédiation s'écrit DANS l'environnement seqExercice, à la suite du
%   \seqCorrige principal :
%
%     \begin{seqExercice}[nom=Calcul mental]
%       Énoncé principal\ldots
%       \seqCorrige{Corrigé principal\ldots}%
%       \seqRemediation{Énoncé de remédiation}{Corrigé de remédiation}%
%     \end{seqExercice}%
%
%   - Le numéro de l'exercice principal (ExoNum) est repris automatiquement
%     pour la remédiation. Pas besoin de propager un argument.
%   - Les exercices SANS \seqRemediation ne produisent rien dans les sections
%     de remédiation (trous autorisés dans la numérotation).
%   - L'énoncé est écrit dans Remediation/re<S>e<E>.tex
%   - Le corrigé est écrit dans Remediation/rc<S>e<E>.tex
%   - Les sections sont compilées en fin de livret par
%     \seqAfficheRemediations puis \seqAfficheCorrigesRemediation.
%
% Patch D4 (v0.11.5) :
%   La macro \seq@ecrit@corrige@serieExos (interne) qui produit le titre
%   « Corrigés des exercices de la série … » écrit désormais un \clearpage
%   au lieu d'un \newpage avant le titre. Cela force la sortie des flottants
%   en attente et garantit qu'une nouvelle série de corrigés démarre sur
%   une page propre, sans superposition possible avec la fin de la série
%   précédente.
% =============================================================================
%
%<*package>
\NeedsTeXFormat{LaTeX2e}%
\ProvidesPackage{seqenseigne-core-exos}%
[2026/05/04 v1.0-dev seqenseigne core exos module]

% =============================================================================
% Dépendances
% =============================================================================
% Tout ce dont on a besoin (multicol, tcolorbox, etoolbox, ifthen, xkeyval)
% est déjà chargé par seqenseigne-core, qui doit être chargé AVANT ce module.
% On ne RequirePackage rien ici pour éviter des doublons.

% =============================================================================
% Configuration de tabularray pour le module
% =============================================================================
% v0.13.5.2 étape 2 :
%   tabularray + bibliothèque booktabs sont déjà chargés par core
%   (\UseTblrLibrary{booktabs}). Les deux \DefTblrTemplate{...}{default}{}%
%   désactivent globalement les templates de caption ("Table N :"
%   au-dessus du tableau) et de capcont pour tous les longtblr du
%   document. Conséquence générale acceptée : plus aucun longtblr
%   n'aura de caption automatique. Note : ceci ne supprime PAS le
%   message "Suite page suivante" qui apparaît quand un tableau long
%   est coupé entre deux pages — celui-ci est géré séparément par
%   tabularray.
\DefTblrTemplate{caption}{default}{}%
\DefTblrTemplate{capcont}{default}{}%

% =============================================================================
% Point d'insertion avant une série d'exercices
% =============================================================================
% Permet d'insérer du contenu entre le titre d'une série et les exercices
\newcommand{\seq@avant@serie}{}%
\newcommand{\seqContenuAvantSerie}[1]{\gdef\seq@avant@serie{#1}}%

% =============================================================================
% Flag d'affichage des corrigés
% =============================================================================
% Initialisé à 'init' ; \seqTitreLivret le passe à 'oui' ou 'non' selon
% l'option [corriges=oui|non]. Pour les séries 0 (révisions) et 1 (fondamentale),
% les corrigés sont écrits SYSTÉMATIQUEMENT (cf. seqSerieExos / seqExercice).
% Pour les autres séries, l'écriture est conditionnelle à \seqCorrigesExos=oui.
\newcommand{\seqCorrigesExos}{init}%

% =============================================================================
% Initialisation des corrigés
% =============================================================================
% Ouvre le fichier CorrList.tex pour accumulation des corrigés au fil de la
% lecture du document. À appeler UNE FOIS, en début de document.
% Note : \seqInitAnnexes (annexes) reste défini dans seqenseigne-core.dtx car
% les annexes sont aussi utilisées par les notions et méthodes (pas seulement
% par les exos).
\newcommand{\seqInitCorriges}{%
	\newcommand{\seq@corrigesoui}{oui}%
	%% ---- Création du fichier pour les corrigés
	\newwrite\corrfile
	\immediate\openout\corrfile=CorrList.tex
}%

% =============================================================================
% v0.11.5 — Initialisation de la remédiation
% =============================================================================
% Ouvre les deux fichiers d'accumulation de la remédiation :
%   RemedEnonceList.tex   — énoncés des exercices de remédiation
%   RemedCorrigeList.tex  — corrigés des exercices de remédiation
% À appeler UNE FOIS, en début de document, en complément de \seqInitCorriges.
%
% Cette macro est OPTIONNELLE : si elle n'est pas appelée, les flux
% \remedenoncefile et \remedcorrigefile ne sont jamais ouverts, et toute
% commande \seqRemediation rencontrée dans le document ne fera rien (un test
% sur \seq@remediation@active évite l'erreur "stream not open").
\newcommand{\seqInitRemediation}{%
	%% Marqueur d'activation, lu par \seqRemediation
	\newcommand{\seq@remediation@active}{oui}%
	%% ---- Création des fichiers de remédiation
	\newwrite\remedenoncefile
	\immediate\openout\remedenoncefile=RemedEnonceList.tex
	\newwrite\remedcorrigefile
	\immediate\openout\remedcorrigefile=RemedCorrigeList.tex
	%% ---- Titres en tête des deux sections de remédiation
	%% On les écrit DÈS l'ouverture des fichiers : ils apparaîtront en tête
	%% lors du \input final, sans mécanisme conditionnel par série (la
	%% remédiation ne distingue pas Révisions/F/A/E — c'est un bloc unique
	%% en fin de livret).
	\immediate\write\remedenoncefile{%
		\noexpand\clearpage
		\noexpand\seqTitreSection{\unexpanded{Exercices de remédiation}}%
	}%
	\immediate\write\remedcorrigefile{%
		\noexpand\clearpage
		\noexpand\seqTitreSectionCorrige{\unexpanded{Corrigés des exercices de remédiation}}%
	}%
}%

% =============================================================================
% Affichage des corrigés en fin de livret
% =============================================================================
% Ferme le fichier CorrList.tex et l'inclut. À appeler UNE FOIS, juste avant
% \end{document}.
\newcommand{\seqAfficheCorriges}{%
	\immediate\closeout\corrfile
	 \clearpage \seqTitreSectionCorrige { Corrigés des exercices de la série exploration } \label {corrigeExos3} 
 \par \nopagebreak \vspace {5pt}\nopagebreak \begin {boiteTitreExo} \par Exercice\;2\par \par \textit {Des caractères sacrés}\end {boiteTitreExo} \par \nopagebreak \vspace {5pt}\nopagebreak   \begin {seqColEnum}[nbCols=1] \item Je sais que $1~302 = 1\times 1~000 + 3 \times 100 + 2 \times 1$. Or, dans le premier nombre en hiéroglyphes, je vois : \begin {seqColItem}[nbCols=1] \item une seule fois le premier hiéroglyphe, \includegraphics [scale=0.2]{N09S01E02_corrige4.png}, donc il vaut 1~000. \item trois fois le deuxième hiéroglyphe, \includegraphics [scale=0.2]{N09S01E02_corrige3.png}, donc il vaut 100. \item deux fois le troisième hiéroglyphe, \includegraphics [scale=0.2]{N09S01E02_corrige1.png}, donc il vaut 1. \end {seqColItem} Je sais que $10~047 = 1 \times 10~000 + 4 \times 10 + 7 \times 1$. Or, dans le deuxième nombre en hiéroglyphes, je vois : \begin {seqColItem}[nbCols=1] \item une seule fois le premier hiéroglyphe, \includegraphics [scale=0.2]{N09S01E02_corrige5.png}, donc il vaut 10~000. \item quatre fois le deuxième hiéroglyphe, \includegraphics [scale=0.2]{N09S01E02_corrige2.png}, donc il vaut 10. \item sept fois le le troisième hiéroglyphe, \includegraphics [scale=0.2]{N09S01E02_corrige1.png} que je connais déjà. \end {seqColItem} Je sais que $2~311~021 = 2 \times 1~000~000 + 3 \times 100~000 + 1 \times 10~000 + 1 \times 1~000 + 2 \times 10 + 1 \times 1$. \newline Or, dans le troisième nombre en hiéroglyphes, je vois : \begin {seqColItem}[nbCols=1] \item deux fois le premier hiéroglyphe, \includegraphics [scale=0.2]{N09S01E02_corrige7.png}, donc il vaut 1~000~000. \item trois fois le deuxième hiéroglyphe, \includegraphics [scale=0.2]{N09S01E02_corrige6.png}, donc il vaut 100~000. \item une seule fois le troisième hiéroglyphe, \includegraphics [scale=0.2]{N09S01E02_corrige5.png} que je connais déjà. \item une seule fois le quatrième hiéroglyphe, \includegraphics [scale=0.2]{N09S01E02_corrige4.png} que je connais déjà. \item deux fois le cinquième hiéroglyphe, \includegraphics [scale=0.2]{N09S01E02_corrige2.png} que je connais déjà. \item une seule fois le sixième hiéroglyphe, \includegraphics [scale=0.2]{N09S01E02_corrige1.png} que je connais déjà. \end {seqColItem} Voici donc dans l'ordre les valeurs des sept hiéroglyphes : \begin {seqColItem}[nbCols=3] \item \includegraphics [scale=0.2]{N09S01E02_corrige1.png} vaut 1. \newline \item \includegraphics [scale=0.2]{N09S01E02_corrige2.png} vaut 10. \newline \item \includegraphics [scale=0.2]{N09S01E02_corrige3.png} vaut 100. \newline \item \includegraphics [scale=0.2]{N09S01E02_corrige4.png} vaut 1~000. \newline \item \includegraphics [scale=0.2]{N09S01E02_corrige5.png} vaut 10~000. \newline \item \includegraphics [scale=0.2]{N09S01E02_corrige6.png} vaut 100~000. \newline \item \includegraphics [scale=0.2]{N09S01E02_corrige7.png} vaut 1~000~000. \end {seqColItem} \item En m'aidant de la question précédente, je peux calculer les nombres pour chaque hiéroglyphe : \newline \includegraphics [scale=0.2]{N09S01E02_enonce4.png} \quad vaut $3 \times 10 + 2 \times 1 = 32$ \newline \includegraphics [scale=0.2]{N09S01E02_enonce5.png} \quad vaut $1 \times 100~000 + 3 \times 10~000 + 2 \times 1 = 130~002$ \newline \includegraphics [scale=0.2]{N09S01E02_enonce6.png} \quad vaut $1 \times 1~000~000 + 2 \times 1~000 + 1 \times 10 = 1~002~010$ \item Je peux écrire $426 = 4 \times 100 + 2 \times 10 + 6 \times 1$ et donc \includegraphics [scale=0.9]{N09S01E02_corrige8.png}. \newline Je peux écrire $527 = 5 \times 100 + 2 \times 10 + 7 \times 1$ et donc \includegraphics [scale=0.9]{N09S01E02_corrige9.png}. \newline Je peux écrire $12~315 = 1 \times 10~000 + 2 \times 1~000 + 3 \times 100 + 1 \times 10 + 5 \times 1$ et donc \includegraphics [scale=0.9]{N09S01E02_corrigeA.png}. \newline Je peux écrire $1~234~000 = 1 \times 1~000~000 + 2 \times 100~000 + 3 \times 10~000 + 4 \times 1~000$ et donc \includegraphics [scale=0.9]{N09S01E02_corrigeB.png}. \newline \end {seqColEnum}  
 
%
}%

% =============================================================================
% v0.11.5 — Affichage des sections de remédiation en fin de livret
% =============================================================================
% Deux macros distinctes : la section « Exercices de remédiation » est typiquement
% placée APRÈS les corrigés principaux, et la section « Corrigés de remédiation »
% en toute fin de livret. Ordre validé v0.11.5 (cadrage Laurent) :
%
%   ... corps du livret ...
%   \seqAfficheAnnexes
%   \seqAfficheCorriges                ← corrigés principaux
%   \seqAfficheRemediations            ← exos de remédiation
%   \seqAfficheCorrigesRemediation     ← corrigés de remédiation
%   ... fiches de résumé éventuelles ...
%
% Si \seqInitRemediation n'a pas été appelée, ces deux macros ne font rien
% (le test sur \seq@remediation@active garantit qu'on n'essaie pas de lire
% un fichier qui n'existe pas).
\newcommand{\seqAfficheRemediations}{%
	\ifcsname seq@remediation@active\endcsname
		\immediate\closeout\remedenoncefile
		 \clearpage \seqTitreSection {Exercices de remédiation} 
%
	\fi
}%
\newcommand{\seqAfficheCorrigesRemediation}{%
	\ifcsname seq@remediation@active\endcsname
		\immediate\closeout\remedcorrigefile
		 \clearpage \seqTitreSectionCorrige {Corrigés des exercices de remédiation} 
%
	\fi
}%

% =============================================================================
% Macro corrigé : écriture dans un fichier séparé via tcbverbatimwrite
% =============================================================================
% \seqCorrige{...} se place DANS un seqExercice. Le contenu est écrit en
% verbatim dans Corriges/c<S>e<E>.tex. L'inclusion dans CorrList.tex se fait
% via \seq@ecrit@corrige@Exo (appelé automatiquement par \begin{seqExercice}).
%
% Avantage du verbatim : permet aux contenus complexes (TikZ, math, etc.)
% de passer sans souci de tokenisation prématurée.
\newcommand{\seqCorrige}[1]{%
	\scantokens{\begin{tcbverbatimwrite}%
			{Corriges/c\theSerieExosNum e\theExoNum.tex} #1
		\end{tcbverbatimwrite}%
	}%
}%

% =============================================================================
% v0.11.5 — Macro de remédiation
% =============================================================================
% Usage :
%   \seqRemediation{énoncé de remédiation}{corrigé de remédiation}%
%
% Se place DANS un seqExercice, après \seqCorrige. Reprend automatiquement
% le numéro courant (ExoNum / SerieExosNum). Ne fait rien si
% \seqInitRemediation n'a pas été appelée — utile pour compiler un livret
% avec ou sans remédiation depuis le même source.
%
% Mécanisme :
%   1. L'énoncé est écrit en verbatim dans Remediation/re<S>e<E>.tex
%   2. Le corrigé est écrit en verbatim dans Remediation/rc<S>e<E>.tex
%   3. Une boîte de titre « Exercice N » est écrite dans RemedEnonceList.tex
%      suivie d'un \input{Remediation/re...tex} (mécanisme parallèle à
%      \seq@ecrit@corrige@Exo)
%   4. Idem pour le corrigé dans RemedCorrigeList.tex
%
% Numérotation :
%   - Le numéro affiché en remédiation est exactement \theExoNum (le numéro
%     principal de l'exercice).
%   - Si certains exos n'ont pas de \seqRemediation, leur numéro est ABSENT
%     de la liste de remédiation (trous autorisés).
\newcommand{\seqRemediation}[2]{%
	\ifcsname seq@remediation@active\endcsname
		%% ---- 1. Écriture verbatim de l'énoncé dans un fichier séparé
		\scantokens{\begin{tcbverbatimwrite}%
				{Remediation/re\theSerieExosNum e\theExoNum.tex} #1
			\end{tcbverbatimwrite}%
		}%
		%% ---- 2. Écriture verbatim du corrigé
		\scantokens{\begin{tcbverbatimwrite}%
				{Remediation/rc\theSerieExosNum e\theExoNum.tex} #2
			\end{tcbverbatimwrite}%
		}%
		%% ---- 3. Entête + \input dans le stream RemedEnonce
		\seq@ecrit@remediation@enonce
		%% ---- 4. Entête + \input dans le stream RemedCorrige
		\seq@ecrit@remediation@corrige
	\fi
}%

% =============================================================================
% Helpers internes : entêtes des exercices de remédiation
% =============================================================================
% Modèle inspiré de \seq@ecrit@corrige@Exo. v0.11.6.3 : la boîte de titre
% suit la même disposition à 2 lignes que le corrigé principal :
%   Ligne 1 : « Exercice N » à gauche, « Objectif XX » à droite si obj=
%   Ligne 2 : titre de l'exercice (italique) si nom=
% L'objectif est passé via l'option [obj=...] de seqExercice ; comme
% l'écriture du flux remédiation se fait au moment où \cmdseq@exercice@obj
% est encore défini (à l'intérieur de seqExercice), la transmission est
% directe — pas de besoin de stocker l'info ailleurs.
\newcommand\seq@ecrit@remediation@enonce{%
	\immediate\write\remedenoncefile{%
		\unexpanded{\par\nopagebreak\vspace{5pt}\nopagebreak%
			\begin{boiteTitreExo}}%
				%% Ligne 1 : numéro à gauche, objectif à droite
				Exercice\noexpand\;\theExoNum
				\ifx\cmdseq@exercice@obj\undefined\else
					\ifnum\value{SerieExosNum}>0
						\unexpanded{\hfill\small\textit}\@charlb Objectif\noexpand\;\eue{\cmdseq@exercice@obj}\@charrb
					\fi
				\fi
				%% Ligne 2 : titre seul
				\ifx\cmdseq@exercice@nom\undefined\else
					\unexpanded{\par\textit}\@charlb \eue{\cmdseq@exercice@nom}\@charrb
				\fi
				\unexpanded{\end{boiteTitreExo}%
			\par\nopagebreak\vspace{5pt}\nopagebreak}%
		%% v0.11.6.2 — \label posé à l'emplacement de l'entête, pour que
		%% la boîte de titre principale (seqExercice avec remediation=oui)
		%% puisse afficher « Remédiation p. \pageref{remed:S-N} ».
		\noexpand\label{remed:\theSerieExosNum-\theExoNum}%
		\noexpand\input{Remediation/re\theSerieExosNum e\theExoNum.tex}%
	}%
}%

\newcommand\seq@ecrit@remediation@corrige{%
	\immediate\write\remedcorrigefile{%
		\unexpanded{\par\nopagebreak\vspace{5pt}\nopagebreak%
			\begin{boiteTitreExo}}%
				%% Ligne 1 : numéro à gauche, objectif à droite
				Exercice\noexpand\;\theExoNum
				\ifx\cmdseq@exercice@obj\undefined\else
					\ifnum\value{SerieExosNum}>0
						\unexpanded{\hfill\small\textit}\@charlb Objectif\noexpand\;\eue{\cmdseq@exercice@obj}\@charrb
					\fi
				\fi
				%% Ligne 2 : titre seul
				\ifx\cmdseq@exercice@nom\undefined\else
					\unexpanded{\par\textit}\@charlb \eue{\cmdseq@exercice@nom}\@charrb
				\fi
				\unexpanded{\end{boiteTitreExo}%
			\par\nopagebreak\vspace{5pt}\nopagebreak}%
		\noexpand\input{Remediation/rc\theSerieExosNum e\theExoNum.tex}%
	}%
}%

% =============================================================================
% v0.11.5 — Cadre de réponse pour l'élève
% =============================================================================
% Usage : \seqCadreReponse{N}  où N est la hauteur du cadre, en lignes.
%
% Produit un cadre tcolorbox avec titre « Réponse » dans les couleurs du
% thème courant (réutilise boiteVariableTitre du module theme). La hauteur
% du cadre est dimensionnée en multiples de \baselineskip pour que N soit
% intuitif côté enseignant ("10 lignes" = 10 \baselineskip de hauteur libre).
%
% Se place typiquement après l'énoncé d'un exercice :
%   \begin{seqExercice}%
%     Calculer 245 + 178.
%     \seqCadreReponse{4}    % 4 lignes pour répondre
%     \seqCorrige{...}%
%   \end{seqExercice}%
%
% Note : le cadre n'apparaît PAS dans les corrigés (il n'est pas écrit
% dans \corrfile).
\newcommand{\seqCadreReponse}[1]{%
	\par\vspace{4pt}%
	\begin{boiteVariableTitre}{Réponse}{\linewidth}%
		\vspace*{#1\baselineskip}%
	\end{boiteVariableTitre}%
	\par\vspace{4pt}%
}%

% =============================================================================
% Helpers internes : entête des séries de corrigés
% =============================================================================
% Écrit dans \corrfile, en début de chaque nouvelle série, le titre de la
% section de corrigés correspondante.
%
% v0.11.5 (patch D4) : \noexpand\newpage → \noexpand\clearpage.
% \newpage termine la page mais laisse les flottants en attente sortir
% sur la suivante ; \clearpage les force à sortir AVANT le saut. Sans
% ce changement, on observait des superpositions entre la fin d'une
% série de corrigés et le début de la suivante quand des flottants
% (figures, tableaux en [h], [t], [b]) étaient en attente.
\newcommand\seq@ecrit@corrige@serieExos{%
	\immediate\write\corrfile{%
		\noexpand\clearpage
		\noexpand\seqTitreSectionCorrige{%
			\ifnum\value{SerieExosNum}=0 \ifnum\value{ExoNum}=0 \unexpanded{Corrigés des exercices de révision} \fi \fi
			\ifnum\value{SerieExosNum}=1 \unexpanded{Corrigés des exercices de la série fondamentale} \fi
			\ifnum\value{SerieExosNum}=2 \unexpanded{Corrigés des exercices de la série avancée} \fi
			\ifnum\value{SerieExosNum}=3 \unexpanded{Corrigés des exercices de la série exploration} \fi
		}%
		\noexpand\label{corrigeExos\theSerieExosNum}%
	}%
}%

\newcommand\seq@ecrit@corrige@Exo{%
	\immediate\write\corrfile{%
		\unexpanded{\par\nopagebreak\vspace{5pt}\nopagebreak%
			\begin{boiteTitreExo}}%
				%% Ligne 1 : numéro à gauche, objectif à droite
				Exercice\noexpand\;\theExoNum
				\ifx\cmdseq@exercice@obj\undefined\else
					\ifnum\value{SerieExosNum}>0
						\unexpanded{\hfill\small\textit}\@charlb Objectif\noexpand\;\eue{\cmdseq@exercice@obj}\@charrb
					\fi
				\fi
				%% Ligne 2 : titre seul
				\ifx\cmdseq@exercice@nom\undefined\else
					\unexpanded{\par\textit}\@charlb \eue{\cmdseq@exercice@nom}\@charrb
				\fi
				\unexpanded{\end{boiteTitreExo}%
			\par\nopagebreak\vspace{5pt}\nopagebreak}%
		\noexpand\input{Corriges/c\theSerieExosNum e\theExoNum.tex}%
	}%
}%

% =============================================================================
% Environnement seqSerieExos : section pour une série d'exercices
% =============================================================================
% Argument : numéro de série (0=révisions, 1=fondamentale, 2=avancée,
% 3=exploration, 4+=série numérotée).
%
% Comportement :
%   - Si on entre dans une nouvelle série (SerieExosNum différent) :
%     met à jour le compteur, RAZ ExoNum, pose un \label
%   - Pour les séries 0 et 1 (révisions/fondamentale), écriture
%     SYSTÉMATIQUE des corrigés
%   - Pour les autres séries, écriture conditionnelle à \seqCorrigesExos=oui
%   - Pose le titre de section (\seqTitreSection) et une boîte d'aide
%     mentionnant la page des corrigés
\newenvironment{seqSerieExos}[1]{%
	\ifnum\value{SerieExosNum}=#1 \else \setcounter{SerieExosNum}{#1}\setcounter{ExoNum}{0} \label{serie:\theSerieExosNum}	\fi
	%% ---- On écrit systématiquement les corrigés pour les révisions et la série fondamentale
	\ifnum\value{SerieExosNum}<2 \seq@ecrit@corrige@serieExos \else \ifdefstring{\seqCorrigesExos}{oui}{ \seq@ecrit@corrige@serieExos }{}	\fi
	\seqTitreSection{%
		\ifnum\value{SerieExosNum}=0 \ifnum\value{ExoNum}=0 Révisions \else Révisions (suite) \fi \fi
		\ifnum\value{SerieExosNum}=1 \ifnum\value{ExoNum}=0 Série fondamentale \else Série fondamentale (suite) \fi \fi
		\ifnum\value{SerieExosNum}=2 \ifnum\value{ExoNum}=0 Série avancée \else Série avancée (suite) \fi \fi
		\ifnum\value{SerieExosNum}=3 \ifnum\value{ExoNum}=0 Série exploration \else Série exploration (suite) \fi \fi
		\ifnum\value{SerieExosNum}>3 \ifnum\value{ExoNum}=0 Série~\theSerieExosNum \else Série~\theSerieExosNum (suite) \fi \fi
	}%
	\seq@avant@serie\gdef\seq@avant@serie{}%
	\setlength{\columnsep}{10pt}%
	\begin{multicols*}{2}%
		\begin{boitePaleBreakNoTitle}%
			\ifnum\value{SerieExosNum}<2%
			\hfill	{\small \textit{Corrigés en page \pageref{corrigeExos\theSerieExosNum}}}%
			\else
			\ifdefstring{\seqCorrigesExos}{oui}{\hfill	{\small \textit{Corrigés en page \pageref{corrigeExos\theSerieExosNum}}}}{}%
			\fi
		}{%
		\end{boitePaleBreakNoTitle}%
	\end{multicols*}%
	\newpage
}%

% =============================================================================
% Environnement seqExercice : un exercice individuel
% =============================================================================
% Options :
%   nom={Titre de l'exercice}        (optionnel)
%   obj={code(s) d'objectif(s)}      (optionnel ; plusieurs séparés par ", ")
%   remediation=oui                  (optionnel ; v0.11.6.2)
%
% Comportement :
%   - Incrémente ExoNum, pose un \label{exo:S-E}%
%   - Pour les séries 0 et 1, écrit le corrigé SYSTÉMATIQUEMENT dans corrfile
%   - Pour les autres séries, écriture conditionnelle à \seqCorrigesExos=oui
%   - Pose une boîte de titre à 3 lignes (v0.11.6.2) :
%       Ligne 1 : « Exercice N » à gauche, « Objectif XX » à droite si obj=
%       Ligne 2 : titre de l'exercice (italique) si nom=
%       Ligne 3 : « Remédiation p. XX » si remediation=oui
%     Les lignes 2 et 3 ne sont émises que si leur contenu est défini.
\define@cmdkey[seq]{exercice}{nom}{}%
\define@cmdkey[seq]{exercice}{obj}{}%
\define@cmdkey[seq]{exercice}{remediation}[non]{}%
\newenvironment{seqExercice}[1][]{%
	%% ---- Options ; remediation=non par défaut, à surcharger par l'utilisateur
	\setkeys[seq]{exercice}{remediation=non,#1}%
	%% ---- Numéro de l'exercice
	\refstepcounter{ExoNum}\label{exo:\theSerieExosNum-\theExoNum}%
	%% ---- On écrit systématiquement les corrigés pour les révisions et la série fondamentale
	\ifnum\value{SerieExosNum}<2 \seq@ecrit@corrige@Exo
	\else \ifdefstring{\seqCorrigesExos}{oui}{ \seq@ecrit@corrige@Exo }{}%
	\fi
	%% ---- Boîte de titre (v0.11.6.2 : 3 lignes)
	\par\nopagebreak\vspace{5pt}\nopagebreak%
	\begin{boiteTitreExo}%
		%% Ligne 1 : numéro à gauche, objectif à droite
		Exercice\;\theExoNum
		\ifx\cmdseq@exercice@obj\undefined\else
			\ifnum\value{SerieExosNum}>0
				\hfill{\small\textit{Objectif\;\cmdseq@exercice@obj}}%
			\fi
		\fi
		%% Ligne 2 : titre seul
		\ifx\cmdseq@exercice@nom\undefined\else
			\par\textit{\cmdseq@exercice@nom}%
		\fi
		%% Ligne 3 : page de la remédiation (v0.11.6.2)
		\ifdefstring{\cmdseq@exercice@remediation}{oui}{%
			\par{\small\textit{Remédiation p.~\pageref{remed:\theSerieExosNum-\theExoNum}}}%
		}{}%
	\end{boiteTitreExo}%
	\par\nopagebreak\vspace{5pt}\nopagebreak%
}{}%

% Fin du module core-exos (partie remediation/exercices)

%% ------------------------------------------------------------------
%% ---- Helpers QCM partagés ----------------------------------------
%% ------------------------------------------------------------------
%% Helpers internes : reconstruction dynamique de l'en-tête de
%% colonnes (A B C ...) du QCM. \seq@qcm@repCodes est reconstruite à
%% chaque ouverture d'environnement seqQcm depuis les compteurs
%% \seq@qcm@iter et \seq@qcm@alphcode (déclarés dans core).
\newcommand\seq@qcm@repCodes{}%

%% ------------------------------------------------------------------
%% ---- \seqQcmQuestion --------------------------------------------
%% ------------------------------------------------------------------
%    \begin{macro}{\seqQcmQuestion}%
%% \seqQcmQuestion{texte de la question}%
%% Produit le numéro de la question (en gras, suivi d'un point), un
%% espace \quad, puis le texte de la question.
%%
%% Le numéro est calculé depuis \therownum (compteur de ligne fourni
%% par tabularray, robuste face aux passes multiples du moteur), en
%% retirant les 2 lignes d'en-tête. Cette approche remplace le
%% compteur \seqQcmItemNum qu'on avait initialement essayé d'utiliser
%% (et qui ne fonctionnait pas, car \stepcounter est appelé à chaque
%% passe de tabularray, faisant sauter la numérotation).
%%
%% Usage côté .tex : à utiliser dans le contexte d'un environnement
%% seqQcm, après un \midrule, avec & et \\ explicites :
%%     \midrule \seqQcmQuestion{Énoncé de la question}%
%%         & rep A & rep B & rep C \\
\newcommand{\seqQcmQuestion}[1]{%
	\textbf{\the\numexpr\therownum-2\relax.}\quad #1
}%
%    \end{macro}%
% -------------------------------------------------------------

%% ------------------------------------------------------------------
%% ---- Environnement seqQcm ----------------------------------------
%% ------------------------------------------------------------------
%    \begin{environment}{seqQcm}%
%% Environnement enfant : produit uniquement le tableau QCM, à
%% utiliser à l'intérieur d'un environnement parent (seqEvalExercice
%% ou seqExercice) qui s'occupe du cartouche, du numéro d'exercice
%% et de l'écriture du corrigé différé.
%%
%% Options de scoring (toutes optionnelles, mais affichage des règles
%% uniquement si les TROIS sont définies — logique tout-ou-rien) :
%%   ptOK, ptPartiel, ptKO : barème pour évaluation. Si l'une des
%%       trois est absente, la dinglist des règles de scoring n'est
%%       pas affichée. Cas typique du QCM hors-éval (révision AE) :
%%       aucune des trois n'est renseignée, pas de scoring affiché.
%% Options de présentation :
%%   nbReps : nombre de colonnes de réponses utilisées effectivement
%%       (de 2 à 5). Le tableau est dimensionné en dur pour 5 colonnes
%%       maximum, l'espace non utilisé reste vide. Défaut : 3.
%%   explication : oui (défaut) / non. Si "non", supprime le bloc
%%       d'explication "Cet exercice est un questionnaire..." pour
%%       laisser à l'utilisateur le soin de le rédiger lui-même.
%%
%% Détection du scoring : on initialise les trois options à \relax,
%% et on teste explicitement si l'utilisateur a défini les trois.
%%
%% Limitations imposées par tabularray :
%%   - colspec figé en \"X[5,c,m]*{5}{X[2,c,m]}\" (pas de macro à
%%     l'intérieur de la spécification de colonnes : tabularray
%%     refuse). D'où la limite à 5 réponses max.
%%   - \cmidrule{2-5} en dur dans l'en-tête (même raison : pas de
%%     macro à l'intérieur du corps de tabularray). Le trait s'adapte
%%     visuellement même avec moins de 5 réponses effectives.
%%   - les \midrule séparateurs entre questions, le \\ de fin de
%%     ligne et le \bottomrule final doivent être écrits par
%%     l'utilisateur, en dehors de toute macro (cf. \seqQcmQuestion
%%     qui ne contient justement aucun de ces éléments structurels).
\define@cmdkey[seq]{seqQcm}{ptOK}{}%
\define@cmdkey[seq]{seqQcm}{ptPartiel}{}%
\define@cmdkey[seq]{seqQcm}{ptKO}{}%
\define@cmdkey[seq]{seqQcm}{nbReps}[3]{}%
\define@choicekey[seq]{seqQcm}{explication}[\val\nr]{oui,non}[oui]{%
	\ifcase\nr\relax
		\def\cmdseq@seqQcm@explication{oui}%
	\or
		\def\cmdseq@seqQcm@explication{non}%
	\fi
}%
\newenvironment{seqQcm}[1][]{%
	%% ---- Réinitialisation des options de scoring (à \relax pour
	%%      détection ultérieure de leur présence)
	\let\cmdseq@seqQcm@ptOK\relax
	\let\cmdseq@seqQcm@ptPartiel\relax
	\let\cmdseq@seqQcm@ptKO\relax
	%% ---- Options
	\setkeys[seq]{seqQcm}{nbReps,explication,#1}%
	%% ---- Bloc d'explication (option explication=oui par défaut)
	\ifx\cmdseq@seqQcm@explication\@undefined\else
		\ifdefstring{\cmdseq@seqQcm@explication}{oui}{%
			Cet exercice est un questionnaire à choix multiples :
			pour chaque question, il y a une ou plusieurs réponse(s) correcte(s). \newline
			Pour répondre, indiquer le numéro de la question et
			recopier la ou les affirmation(s) juste(s) sur votre copie. \newline
			\textbf{Aucune justification n'est attendue pour ce QCM.} \newline
		}{}%
	\fi
	%% ---- Bloc des règles de scoring (uniquement si les TROIS
	%%      paramètres ptOK/ptPartiel/ptKO sont définis)
	\ifx\cmdseq@seqQcm@ptOK\relax\else
		\ifx\cmdseq@seqQcm@ptPartiel\relax\else
			\ifx\cmdseq@seqQcm@ptKO\relax\else
				Pour chaque question:
				\begin{dinglist}{224}%
					\item Donner toutes les bonnes réponses rapporte \cmdseq@seqQcm@ptOK.
					\item Ne pas donner toutes les bonnes réponses, mais sans mauvaise réponse,
						rapporte \cmdseq@seqQcm@ptPartiel.
					\item Ne pas donner de réponse n'est pas pénalisé.
					\item Donner au moins une mauvaise réponse enlève \cmdseq@seqQcm@ptKO.
				\end{dinglist}%
			\fi
		\fi
	\fi
	%% ---- Construction de la ligne d'en-tête (A B C...)
	\setcounter{seq@qcm@alphcode}{40}%
	\renewcommand\seq@qcm@repCodes{}%
	\forloop{seq@qcm@iter}{0}{\value{seq@qcm@iter} < \cmdseq@seqQcm@nbReps}%
	{\stepcounter{seq@qcm@alphcode} \eappto\seq@qcm@repCodes{ & \char"\theseq@qcm@alphcode}}%
	%% ---- Le tableau lui-même
	%% colspec figé en X[5,c,m]*{5}{X[2,c,m]} : tabularray refuse
	%%   les macros dans la spécification de colonnes. Cohérence
	%%   visuelle même si nbReps < 5.
	%% rowhead=2 : les 2 premières lignes (Questions/Affirmations
	%%   puis A B C) sont répétées en tête de chaque page.
	%% expand=\seq@qcm@repCodes : injecte les codes A B C qui ont été
	%%   construits dynamiquement plus haut.
	\begin{center}%
	\begin{longtblr}[expand=\seq@qcm@repCodes]{%
			width = 0.9\linewidth,
			colspec = {X[5,c,m]*{5}{X[2,c,m]}},
			rowhead = 2,
			rowsep = 2pt
		}%
		\toprule \textbf{Questions} & \textbf{Affirmations} \\
		\cmidrule{2-5} \seq@qcm@repCodes \\
		}{%
	\end{longtblr}%
	\end{center}%
}%
%    \end{environment}%
% -------------------------------------------------------------

% Fin du module core-exos
%</package>
%
% \endinput
